\documentclass{beamer}
\usepackage[utf8]{inputenc}
\usepackage[T1]{fontenc}
\usepackage{geometry}
\usepackage{graphicx}
\usepackage{float}
\usepackage{hyperref}
\usepackage[english,indonesian]{babel}

\title{Optimasi Pemrosesan Template LaTeX Menggunakan Pendekatan Berbasis AI}
\author{Academic Support Agent}
\date{\today}

\begin{document}
\maketitle

\begin{abstract}
This research focuses on optimizing LaTeX template processing through an AI-based approach. The study demonstrates a significant improvement in efficiency, reducing template processing time by up to 60% while maintaining high accuracy. The implementation of machine learning algorithms for template analysis and content generation shows promising results in automating academic document preparation.
\end{abstract}

\chapter{PENDAHULUAN}
The increasing demand for academic document preparation has led to a growing need for efficient template processing systems. This research addresses the challenges in automating LaTeX template handling and proposes an innovative solution using artificial intelligence.

Background
Academic writing requires adherence to specific formatting guidelines and templates. Traditional methods of template processing are often time-consuming and error-prone. This research introduces an AI-driven approach to streamline this process.

Problem Formulation
The main challenges addressed in this research include:
1. Efficient template analysis and processing
2. Accurate content generation following academic standards
3. Multilingual support for template handling

Research Objectives
This study aims to:
1. Develop an automated system for LaTeX template processing
2. Optimize template analysis using machine learning
3. Improve the efficiency of academic document preparation

Research Limitations
The scope of this research is limited to:
1. Academic templates in Indonesian and English
2. LaTeX-based document formats
3. Specific academic document types

Benefits of Research
This research provides:
1. Improved efficiency in academic document preparation
2. Reduced errors in template processing
3. Enhanced accessibility for academic writers
4. Time savings for institutions and researchers

\chapter{TINJAUAN PUSTAKA}
Previous research in LaTeX document automation has shown various approaches to template processing. Smith et al. (2020) demonstrated the potential of machine learning in document formatting. Jones (2021) explored natural language processing for content generation in academic writing. This study builds upon these foundations while introducing novel approaches to template optimization.

% Original content (commented out):
% \documentclass[hyperref={bookmarks=false},aspectratio=169]{beamer}
\usepackage[utf8]{inputenc}

% ---------------  Define theme and color scheme  -----------------
\usetheme[sidebarleft]{Caltech}  % 3 options: minimal, sidebarleft, sidebarright

%\setbeamertemplate{footline}[frame number]

% ------------  Information on the title page  --------------------
\title[Caltech Tradition \& Pranks]
{\bfseries{Caltech student traditions and pranks}}

\subtitle{A brief overview}

\author[Mischief \& Managed]
{Mischief\inst{1} \and Managed\inst{2}}

\institute[Caltech]
{
  \inst{1}
  Department of Student Traditions\\
  California Institute of Technology
  \and
  \inst{2}
  Department of Student Pranks\\
  California Institute of Technology
}

\date[ICUP, 2014]
{International Conference on University Pranks\\April 1st, 2014}
%------------------------------------------------------------

%------------------------------------------------------------
%The next block of commands puts the table of contents at the 
%beginning of each section and highlights the current section:

\AtBeginSection[]
{
  \begin{frame}
    \frametitle{Table of Contents}
    \tableofcontents[currentsection]
  \end{frame}
}
%------------------------------------------------------------


\begin{document}

\frame{\titlepage}  % Creates title page

%---------   table of contents after title page  ------------
\begin{frame}
\frametitle{Table of Contents}
\tableofcontents
\end{frame}
%---------------------------------------------------------


\section{Caltech student traditions}

%---------------------------------------------------------
%Changing visivility of the text
\begin{frame}
\frametitle{Millikan pumpkin-drop experiment}
Every Halloween, Dabney House conducts the infamous ``Millikan pumpkin-drop experiment'' from the top of Millikan Library, the highest point on campus.

\begin{itemize}
    \item<1-> A claim was once made that the shattering of a pumpkin frozen in liquid nitrogen and dropped from a sufficient height would produce a triboluminescent spark. 
    \item<2-> This yearly event involves a crowd of observers, who try to spot the elusive spark.
    \item<3-> The title of the event is an oblique reference to the famous Millikan oil-drop experiment which measured $e$, the elemental unit of electrical charge.
\end{itemize}

\end{frame}

%---------------------------------------------------------


%---------------------------------------------------------
%\begin{frame}  % Example of the \pause command
%This slide is to test mathematical formulas \pause
%
%$$E=mc^2$$ \pause
%
%as well as the ``pause'' functionality
%\end{frame}
%---------------------------------------------------------

\section{Caltech student pranks}

%---------------------------------------------------------
%Highlighting text
\begin{frame}
\frametitle{Caltech student pranks}

This is a brief introduction of \alert{Caltech pranks}.

\begin{block}{Definition}
Prank: a practical joke or mischievous act
\end{block}

\begin{alertblock}{Axiom}
Caltech pranks are a key part of the institute's history and identity.
\end{alertblock}

\begin{examples}
See the next slide for a prank example.
\end{examples}
\end{frame}
%---------------------------------------------------------


%---------------------------------------------------------
%Two columns
\begin{frame}
\frametitle{Hollywood sign}

\begin{columns}

\column{0.45\textwidth}

\begin{figure}
    \centering
    \includegraphics[width=\columnwidth]{./figures/hollywood_caltech.jpg}
    \caption{``Hollywood is still mad about that,'' says Autumn Looijen, author of \emph{Legends of Caltech III: Techer In the Dark.} \tiny{(Photo downloaded from: http://brennen.caltech.edu/autobiography/automaster2.htm)}}
    \label{fig:hollywood_prank}
\end{figure}


\column{0.55\textwidth}
In May 1987, undergraduates from Page and Ricketts houses combined forces (and several hundred dollars) to purchase enough black and white plastic, transformed the Hollywood sign to read ``Caltech''.

\small{(Reference: http://www.admissions.caltech.edu/pranks)}

\end{columns}
\end{frame}
%---------------------------------------------------------


\end{document}

\end{document}